\section{Contraction at sparse representation levels}\label{sec:sparse}
\label{ac:part:51}

We now contract an irreducible that first appears on a short list of cards,
even after a comparable number of trajectories has been prescribed. The row
case uses an alternating placement kernel. The column case uses the same
placement representation twisted by sign.

\begin{lemma}[Conditional sparse contraction]\label{lem:sparse}
For every \(0<\theta<1\) and \(K_0\ge1\), there are \(\rho>0\) and \(R_0\)
such that, for each power of two \(R\ge R_0\), one can choose
\(z_{\rm sp}(R)>0\) with the following property. For every allowed grid and
prescribed family \(\mathcal H\), suppose that the integer \(k\ge1\) satisfies
\begin{equation}
 h\le K_0 k,\qquad h+k\le s^{1-\theta},
 \label{ac:eq:5}
\end{equation}
and let \(\lambda\vdash M=s-h\) satisfy
\(\lambda_1=M-k\) or \(\lambda'_1=M-k\). Then, for every
\(z\in[0,z_{\rm sp}(R)]\),
\begin{equation}
 \|K_\lambda^{\mathcal H}\|_{\rm op}\le s^{-\rho k}.
 \label{ac:eq:6}
\end{equation}
\end{lemma}
\begin{proof}
Put \(B_0=\lceil8/\theta\rceil\). For \(b>B_0\) we use only the line comparison
\eqref{ac:eq:3}. For \(b\le B_0\) we first prove the estimate at \(z=0\), then
use continuity after fixing \(R\). The implicit constants in the entry
estimates below are absolute; their absorption into powers of \(s\) will
determine how large \(R\) must be in terms of \(\theta,K_0\).

First assume \(\lambda_1=M-k\). For input and output \(k\)-placements \(x,y\),
let \(p_A(x_A,y_A)\) be the conditional probability of the paths for the
labels in \(A\subseteq[k]\), with \(p_\varnothing=1\). Define a kernel from
input placements to output placements by
\begin{equation}
 Q(x,y)=s^{-k}\sum_{A\subseteq[k]}(-1)^{k-|A|}s^{|A|}p_A(x_A,y_A).
 \label{ac:eq:7}
\end{equation}
Its \(A=[k]\) term is the full placement transition kernel. For a proper
subset \(A\), the corresponding operator has output range among the
pullbacks \(y\mapsto f(y_A)\) of functions on \(|A|\)-placements. This
subrepresentation contains only types with first row at least
\(M-|A|>M-k\), by the branching rule in Section~\ref{sec:prelim}. Projection
to type \(\lambda\) therefore kills every term for a proper subset. The full
placement representation contains \([\lambda]\), so its surviving block is a
copy of \(K_\lambda^{\mathcal H}\). Consequently
\(\|K_\lambda^{\mathcal H}\|_{\rm op}\le\|Q\|_{\rm HS}\).
The scale \(s\) in \eqref{ac:eq:7} will cancel the probability \(1/s\) of a
path whose stage lines are untouched by every other constraint.

If the paths of \(A\) augment \(\mathcal H\) to a disjoint family
\(\mathcal H_A\), then \eqref{ac:eq:4} and
\(C(\mathcal H_A)\le b(h+|A|)\) give a crude bound. Invalid subsets have
probability zero. Summing over \(A\) yields
\begin{equation}
 |s^kQ(x,y)|\le 2^k\exp\{bh+(1+\log4)bk\}.
 \label{ac:eq:8}
\end{equation}

To count these entries, extend \(Q\) by zero when either endpoint array fails
to be an injection into its available input or output slots. On endpoint
placements we keep the alternating formula \eqref{ac:eq:7}, even when the
full family of \(k\) paths collides at an intermediate layer. Now choose all
input and output slots independently and uniformly from the full grid
\(\Omega\), with replacement. The zero extension gives the exact identity
\[
 \|Q\|_{\rm HS}^2
 =s^{-2k}\sum_{x,y\in\Omega^k}|s^kQ(x,y)|^2
 =\mathbb E_{\rm full}|s^kQ(x,y)|^2.
\]
Every sampled pair of endpoints determines an auxiliary trajectory by the
layer formula in Section~\ref{sec:paths}. At any fixed layer its slot is
uniform in \(\Omega\), and the slots for different particles are independent.
The following counting estimates use this independence between particles;
the conditional shuffle probabilities remain inside \(Q\).

We will count two kinds of sharing: a common slot at a layer, or a common line
at a stage. Suppose \(\omega\) bounds the probability that one independent
particle shares with any fixed trajectory. Form the graph whose vertices are
the particles and holes, with an edge for each sharing involving a particle.
For \(0\le w\le k\), if at least \(w\) particles are incident to edges, this graph contains a
directed forest with at least \(\lceil w/2\rceil\) edges and distinct particle
tails. To see this, in a component meeting holes, point each particle to a
neighbor one step closer to the nearest hole. Distances decrease along these
edges, so they form a forest rooted at holes, with every particle a tail. In
a component containing only particles, with \(v\ge2\) incident vertices, a rooted spanning
tree supplies \(v-1\ge v/2\) particle tails.

For a fixed directed forest with \(j\) edges, condition on all paths except a
leaf particle. Its parent is now fixed, its sharing event has probability at
most \(\omega\), and it occurs in no other remaining edge. Removing leaves
one by one bounds the forest probability by \(\omega^j\). There are at most
\(\binom{k}{j}(k+h)^j\) choices of its tails and parents. Thus, if
\(r=(k+h)\omega\le1\), the probability that at least \(w\) particles are
incident to sharing is at most
\begin{equation}
 \sum_{j=\lceil w/2\rceil}^{k}\binom{k}{j}r^j
 \le 2^k r^{w/2}.
 \label{ac:eq:9}
\end{equation}

For \(b>B_0\), take sharing to mean a common stage line. A particle uses a
fixed line at stage \(j\) with probability \(m_j/s\), so
\(\omega\le b\max_jm_j/s\). Since \(\max_jm_j\le s^{2/b}\),
\[
 (k+h)\omega\le b\,s^{-\theta+2/b}\le s^{-\theta/2}
\]
for sufficiently large \(R\), uniformly in \(b>B_0\). Here \(2/b<\theta/4\),
and \(b\le s^{\theta/4}\) follows uniformly from \(s\ge R^b\) after increasing
\(R\).

If a path shares no stage line with another particle path or a hole, adding
it to any subset multiplies that subset's transition probability by \(1/s\).
Its lines are independent of all the other constraints, and each uses the
uniform marginal for one slot. The statement also holds when the subset has
probability zero. The paired terms in \eqref{ac:eq:7} therefore cancel.
Thus \(Q\ne0\) requires all \(k\) particles to be incident to line sharing.
Using \eqref{ac:eq:9} with \(w=k\) and then \eqref{ac:eq:8},
\[
\begin{aligned}
 \Pr_{\rm full}(Q\ne0)&\le2^k s^{-\theta k/4},\\
 \|Q\|_{\rm HS}
 &\le2^{3k/2}e^{bh+(1+\log4)bk}s^{-\theta k/8}
 \le s^{-\theta k/16}
\end{aligned}
\]
for large \(R\). The last step uses \(h\le K_0k\) and
\(b/\log s\le1/\log R\).

\subsection{A bounded number of coordinates}
\label{ac:part:85}
When \(b\le B_0\), a line can occupy a substantial fraction of the grid, so
many paths may share a line. We work at \(z=0\) and use the fact that most
paths still see only a small fraction of each line prescribed. Put
\[
 \xi=\frac1{10B_0},\qquad \eta=s^{-\xi},
\]
and increase \(R\) so that \(\eta\le1/2\) for all \(s\ge R\).

For any endpoint arrays, let \(t_L\) be the number of their \(k\) auxiliary
paths using line \(L\), whether or not those paths are jointly feasible. Call
\(L\) light if \(h_L+t_L\le |L|\eta\). A particle is good if every line on its
path is light and it shares no slot at any layer with another particle or a
hole. Let \(G\) be the set of good particles. We first prove the entry bound
\begin{equation}
 |s^kQ|\le \exp(O(b(k+h)))\,\eta^{|G|/2}.
 \label{ac:eq:10}
\end{equation}
It is immediate for the zero extension outside endpoint placements, so
consider an entry of the original placement matrix. For a subset \(A\) whose
paths are compatible with \(\mathcal H\) at every layer, the uniform
conditional line law gives
\[
 s^{|A|}p_A(x_A,y_A)
 =\prod_L\frac{m^{v_L(A)}}{(m-h_L)_{v_L(A)}},
 \qquad m=|L|,
\]
where \(v_L(A)\) counts its paths on \(L\). For an invalid subset \(A\), its
probability is zero and this product is not formed.

Because a good path shares no layer slot with another path or a hole,
\(J\cup S\) is valid exactly when \(J\) is valid for every
\(J\subseteq[k]\setminus G\) and \(S\subseteq G\). We may therefore group the
alternating sum by \(J\) and discard every invalid \(J\) before forming line
factors. For a valid \(J\), let \(r_L(J)\) count its paths on \(L\), and
define
\[
 D_J=\prod_{L\text{ non-light}}
       \frac{m^{r_L(J)}}{(m-h_L)_{r_L(J)}},\qquad m=|L|.
\]
All these denominators are positive. The bound for \(\phi_m\) in
Section~\ref{sec:paths} gives
\[
\begin{aligned}
 \log D_J
 &=\sum_{L\text{ non-light}}
    \bigl(\phi_{|L|}(h_L+r_L(J))-\phi_{|L|}(h_L)\bigr)
 \\
 &\le\sum_L\phi_{|L|}(h_L+r_L(J))
 \le b(h+|J|).
\end{aligned}
\]

At each light line introduce \(W_L=m/Z_L\), where the \(Z_L\) are independent
gamma variables with shape \(m-h_L+1\) and unit rate. The elementary identity
\[
 \mathbb E Z^{-v}=\frac{\Gamma(\alpha-v)}{\Gamma(\alpha)}
 \quad\text{for }Z\sim\operatorname{Gamma}(\alpha,1),\quad 0\le v<\alpha,
\]
follows by integrating the gamma density. It gives
\(\mathbb E W_L^v=m^v/(m-h_L)_v\) for the exponents used here. Lightness makes
these moments finite. If \(L_j(i)\) denotes the stage-\(j\) line of a good
path \(i\), alternation over the good paths now gives the exact expression
\[
\begin{split}
 s^kQ(x,y)
 =\sum_{\substack{J\subseteq[k]\setminus G\\J\text{ valid}}}
 &(-1)^{k-|J|-|G|}D_J\\
 &\times\mathbb E\!\left[
   \prod_{L\text{ light}}W_L^{r_L(J)}
   \prod_{i\in G}\left(\prod_{j=1}^bW_{L_j(i)}-1\right)\right].
\end{split}
\]
Every path using a light line uses the same variable \(W_L\). We will group all
factors by line before taking moments; independence is available between
lines, not between particle factors.

For nonnegative integers \(u,v\) with \(u+v\le t_L\), density tilting gives
\begin{equation}
 \mathbb E W_L^v|W_L-1|^u
 =\frac{m^{u+v}}{(m-h_L)_{u+v}}\,
       \mathbb E|1-Z'/m|^u
 \le C_1^{u+v}\eta^{u/2},
 \label{ac:eq:11}
\end{equation}
where \(Z'\) has gamma shape \(n=m-h_L+1-u-v\) and unit rate, and \(C_1\) is absolute.
Indeed lightness gives the margin
\[
 n\ge m(1-\eta)+1>0,\qquad
 \frac{m^{u+v}}{(m-h_L)_{u+v}}
 \le(1-\eta)^{-(u+v)}\le2^{u+v}.
\]
For a gamma variable of shape \(n\), the centered moment generating function
is \(e^{-nt}(1-t)^{-n}\); its logarithm is at most \(Cnt^2\) for
\(|t|\le1/2\). Chernoff's inequality gives
\(\Pr(|Z'-n|>x)\le2\exp[-c\min\{x^2/n,x\}]\). Integrating this tail yields
\[
 (\mathbb E|Z'-n|^u)^{1/u}\le C(\sqrt{nu}+u),\qquad u\ge1.
\]
When \(u\ge1\), we have \(u/m\le\eta\), \(n\le m+1\), and
\(\lvert1-n/m\rvert\le\eta+1/m\le2\eta\), since \(1\le u\le m\eta\).
The triangle inequality therefore gives
\[
 (\mathbb E|1-Z'/m|^u)^{1/u}
 \le |1-n/m|+\frac C m(\sqrt{nu}+u)
 \le C'\sqrt\eta.
\]
This proves \eqref{ac:eq:11}; when \(u=0\), only the prefactor bound is needed.

For each good path, telescope in stage order:
\[
 \prod_{j=1}^bW_{L_j(i)}-1
 =\sum_{a=1}^b (W_{L_a(i)}-1)\prod_{j<a}W_{L_j(i)}.
\]
In one expanded term, let \(u_L\) count its factors \(W_L-1\), and \(v_L\)
its undifferenced factors \(W_L\). Each good path supplies one difference,
so \(\sum_Lu_L=|G|\). A tracked path supplies at most one factor at any
line at a given stage, so \(u_L+v_L\le t_L\). Independence across lines and
\eqref{ac:eq:11} bound the expectation of that term in absolute value by
\(C_1^{bk}\eta^{|G|/2}\). There are at most \(b^{|G|}\) telescoping choices
and \(2^{k-|G|}\) choices of valid \(J\). Together with the bound on \(D_J\),
these prove \eqref{ac:eq:10}.

It remains to count good particles. All probabilities in this count again
refer to the independent endpoint arrays sampled with replacement in the full grid;
the holes are fixed. The sparse hypothesis and
\(\xi=1/(10B_0)\le\theta/80\) give, for large \(R\),
\[
 \frac{(b+1)(k+h)}s\le s^{-\theta/2},\qquad
 \frac{2h}{s\eta},\ \frac{2k}{s\eta}
 \le2s^{-(\theta-\xi)}\le s^{-\theta/2}.
\]
In particular these quantities are at most one.

If \(|G|<k/2\), either at least \(k/4\) particles share a layer slot with
another path, or at least \(k/4\) particles use a non-light line. The first
event is bounded by \eqref{ac:eq:9} with \(\omega=(b+1)/s\). In the second
event, some stage \(j\) has at least \(k/(4B_0)\) particles on non-light
lines. Each such line satisfies \(h_L>m_j\eta/2\) or \(t_L>m_j\eta/2\), so
one of these two classes receives at least
\(a_k=k/(8B_0)\) particles.

At a fixed stage, the lines with \(h_L>m_j\eta/2\) are determined by the
holes and occupy a fraction at most \(2h/(s\eta)\) of the slots. The
independent occupancy bound for at least \(a_k\) particles in this fixed set
is \(2^k(2h/(s\eta))^{a_k}\). For the lines with \(t_L>m_j\eta/2\), their
number is at most \(r_j=\lfloor2k/(m_j\eta)\rfloor\). If \(r_j=0\), there
are none. Otherwise fix a set of at most \(r_j\) lines first. A particle
uses it with probability at most \(r_jm_j/s\le2k/(s\eta)\), so its occupancy
bound is \(2^k(2k/(s\eta))^{a_k}\). We then sum over the at most
\((s+1)^{r_j}\) possible fixed sets. Combining the three events gives
\[
\begin{split}
 \Pr_{\rm full}(|G|<k/2)\le {}&
 2^k\left(\frac{(b+1)(k+h)}s\right)^{k/8}
 +b\,2^k\left(\frac{2h}{s\eta}\right)^{a_k}\\
 &+2^k\sum_{\substack{1\le j\le b\\r_j\ge1}}
       (s+1)^{r_j}\left(\frac{2k}{s\eta}\right)^{a_k}.
\end{split}
\]
The thresholds may be nonintegral: taking their ceilings only decreases
these powers of bases in \([0,1]\).

The cost of choosing the lines with \(t_L>m_j\eta/2\) is uniformly small. Since
\(b\le B_0\) and \(R\le m_j\le R^2\),
\[
 m_j\ge s^{1/(2B_0)},\qquad m_j\eta\ge s^{2/(5B_0)},\qquad
 \frac{r_j\log(s+1)}{k\log s}
 \le2s^{-2/(5B_0)}\frac{\log(s+1)}{\log s}=o_R(1).
\]
Before the factors \(2^k\), the stage sum, and this set count, the vertex
term saves \((\theta/16)k\log s\) and each term for non-light lines saves
\((\theta/(16B_0))k\log s\). The displayed ratio tends to zero uniformly for
\(s\ge R\); the other two factors also cost \(o_R(k\log s)\). Increasing
\(R\) therefore gives
\begin{equation}
 \Pr_{\rm full}(|G|<k/2)\le s^{-\nu k},
 \qquad \nu=\frac{\theta}{100B_0}>0,
 \label{ac:eq:12}
\end{equation}
uniformly over the grids with \(b\le B_0\). On \(|G|\ge k/2\),
\eqref{ac:eq:10} gives the entry saving \(s^{-\xi k/4}\). On the complement,
use \eqref{ac:eq:8} and the square root of \eqref{ac:eq:12} in the
Hilbert--Schmidt expectation. We obtain
\[
 \|Q\|_{\rm HS}
 \le \exp(O(b(k+h)))
       \left(s^{-\xi k/4}+s^{-\nu k/2}\right).
\]

\subsection{Column types and the common perturbation interval}

Suppose now that \(\lambda'_1=M-k\). The sign twist of the \(k\)-placement
representation contains \([\lambda]\). Its entry for placements \(x,y\) is
the conditional average of
\(\operatorname{sgn}(g)\mathbf1_{\{g(x)=y\}}\), where \(g\) is the remaining
bijection in fixed slot identifications. An entry whose full path family is
invalid is zero. For a valid entry, condition on all \(h+k\) paths.

On each line, list the assigned paths first in the same order at input and
output, followed by the unassigned sites. Deleting the fixed assignments
then leaves a residual bijection, and comparison with the original slot
orders contributes only a deterministic sign. Applying this to the lines and
composing the stages expresses the sign of the remaining bijection as a fixed
factor times the product of the residual line signs. The conditioned lines
are independent. It is therefore enough to find one line whose residual sign
has mean zero.

For \(b>B_0\),
\[
 \frac{s}{\max_jm_j}\ge s^{1-2/b}>s^{1-\theta}\ge h+k.
\]
Every stage has more lines than paths, so it has a line untouched by all
\(h+k\) paths. Its residual law is \(L_m(z)\), which has expected sign zero.
For \(b\le B_0\) at \(z=0\), fix any stage \(j\). The number of unassigned
sites satisfies
\[
 s-h-k\ge s(1-s^{-\theta})>\frac{s}{m_j}
\]
once \(R^{-\theta}+R^{-1}<1\). There are \(s/m_j\) lines in that stage, so one
line has at least two unassigned sites. Its residual bijection under the
uniform conditional law is uniform between the unassigned site sets and has expected sign
zero. Thus every signed placement entry vanishes in the relevant parameter
range of each branch.

We finish by fixing the constants in the promised order. Set
\(\beta=\min\{\xi/4,\nu/2\}>0\). For an absolute constant \(C\), the
bounded-coordinate row estimate is at most
\[
 2e^{Cb(1+K_0)k}s^{-\beta k}\le s^{-\beta k/2}
\]
for all sufficiently large \(R\), since \(b/\log s\le1/\log R\).
The corresponding
column operator is zero at \(z=0\). The many-coordinate row estimate is
at most \(s^{-\theta k/16}\), and its column operator is zero whenever
\eqref{ac:eq:3} holds. We may therefore choose
\[
 \rho=\min\{\theta/32,\beta/4\}>0
\]
before \(R\), then choose \(R_0\) large enough for all the preceding estimates.
For every \(R\ge R_0\), the bounded-coordinate bounds at \(z=0\) have strict
slack relative to \(s^{-\rho k}\).

Fix now a power of two \(R\ge R_0\). When \(b\le B_0\), there are finitely
many allowed grids, with \(s\le R^{2B_0}\), and finitely many feasible path
families, integers \(k\), and partitions. Each conditional operator is a
finite expression in line probabilities whose conditioning denominators are
positive at \(z=0\); it is therefore continuous there. Strict slack and
finiteness give \(\delta(R)>0\) such that \eqref{ac:eq:6} holds for all these
bounded-coordinate operators when \(0\le z\le\delta(R)\). Set
\[
 z_{\rm sp}(R)=\min\{\delta(R),z_{\rm cmp}(R)\}>0.
\]
The many-coordinate proof applies throughout this interval by
\eqref{ac:eq:3}, so it covers every \(b\) and completes the lemma.
\end{proof}
